\documentclass[11pt,a4paper]{article}
% ==============================================================
%  Préambule commun aux fiches -- Advanced Probabilities (M1)
%  Usage, dans chaque fiche :
%     \documentclass[11pt,a4paper]{article}
%     % ==============================================================
%  Préambule commun aux fiches -- Advanced Probabilities (M1)
%  Usage, dans chaque fiche :
%     \documentclass[11pt,a4paper]{article}
%     % ==============================================================
%  Préambule commun aux fiches -- Advanced Probabilities (M1)
%  Usage, dans chaque fiche :
%     \documentclass[11pt,a4paper]{article}
%     \input{../preambule.tex}
%     \fiche{C5}{Vecteurs gaussiens}{Gaussian vectors}
% ==============================================================
\usepackage[T1]{fontenc}
\usepackage[utf8]{inputenc}
\usepackage[french]{babel}
\usepackage{lmodern}
\usepackage[a4paper,margin=1.9cm,top=1.7cm,bottom=1.9cm]{geometry}
\usepackage{amsmath,amssymb,amsthm,mathtools}
\usepackage[inline]{enumitem}
\usepackage{booktabs,array,tabularx}
\usepackage{xcolor}
\usepackage[most]{tcolorbox}
\usepackage{fancyhdr}
\usepackage{hyperref}
\hypersetup{colorlinks=true,linkcolor=blue!50!black,urlcolor=blue!50!black}

\setlist{itemsep=1pt,topsep=3pt,parsep=0pt}
\setlength{\parindent}{0pt}
\setlength{\parskip}{3pt}
\emergencystretch=3em

% ---------- Identification de la fiche -------------------------
\newcommand{\ficheid}{}
\newcommand{\fichetitre}{}
\newcommand{\fiche}[3]{%
  \renewcommand{\ficheid}{#1}\renewcommand{\fichetitre}{#2}%
  \pagestyle{fancy}\fancyhf{}%
  \fancyhead[L]{\small\textsf{Advanced Probabilities -- Fiche #1}}%
  \fancyhead[R]{\small\textsf{#2}}%
  \fancyfoot[C]{\small\thepage}%
  \begin{center}
    {\LARGE\bfseries Fiche #1 -- #2}\par\smallskip
    {\large\itshape (#3)}\par\smallskip
    {\small Cours de V.\,M.~Lambert, M1 -- semestre 1}
  \end{center}\vspace{-2pt}\hrule\vspace{6pt}}

% ---------- Environnements numérotés (un seul compteur) --------
\theoremstyle{definition}
\newtheorem{definition}{Définition}
\renewcommand{\thedefinition}{\ficheid.\arabic{definition}}
\newtheorem{theoreme}[definition]{Théorème}
\newtheorem{proposition}[definition]{Proposition}
\newtheorem{lemme}[definition]{Lemme}
\newtheorem{corollaire}[definition]{Corollaire}
\newtheorem{exemple}[definition]{Exemple}
\theoremstyle{remark}
\newtheorem{remarque}[definition]{Remarque}

% Résultat hors cours (standard, signalé comme tel)
\theoremstyle{definition}
\newtheorem{horscours}[definition]{Hors cours}

% Idée de preuve
\newenvironment{idee}{\par\small\textit{Idée de preuve.}~}{\par}

% ---------- Encadrés -----------------------------------------
\tcbset{fichebox/.style={enhanced,breakable,boxrule=0.6pt,arc=2pt,
  left=5pt,right=5pt,top=3pt,bottom=3pt,fonttitle=\bfseries}}
\newtcolorbox{piege}[1][]{fichebox,colback=red!4,colframe=red!60!black,
  title={Pièges classiques},#1}
\newtcolorbox{liens}[1][]{fichebox,colback=blue!3,colframe=blue!50!black,
  title={Liens entre notions},#1}
\newtcolorbox{methode}[1][Méthode]{fichebox,colback=green!4,colframe=green!40!black,title={#1}}
\newtcolorbox{alerte}[1][]{fichebox,colback=orange!6,colframe=orange!70!black,
  title={Attention à l'énoncé},#1}

% Exercice de TD (fiches de rappels)
\newcounter{exo}
\newenvironment{exo}[1][]{\refstepcounter{exo}\par\medskip
  \noindent{\large\bfseries Exercice~\theexo}\ifx\relax#1\relax\else~{\large\bfseries -- #1}\fi\par\nopagebreak\smallskip}{\par}
\newcommand{\outils}{\textbf{Outils du cours.}~}
\newcommand{\meth}{\textbf{Méthode / astuce.}~}
\newcommand{\indic}{\textbf{Indication.}~}
\newcommand{\rep}{\textbf{Réponse.}~}

% ---------- Notations du cours ---------------------------------
\newcommand{\R}{\mathbb{R}}
\newcommand{\N}{\mathbb{N}}
\newcommand{\Q}{\mathbb{Q}}
\newcommand{\C}{\mathbb{C}}
\newcommand{\PP}{\mathbb{P}}
\newcommand{\QQ}{\mathbb{Q}}
\newcommand{\E}{\mathbb{E}}
\newcommand{\F}{\mathcal{F}}
\newcommand{\G}{\mathcal{G}}
\newcommand{\Hc}{\mathcal{H}}
\newcommand{\Tc}{\mathcal{T}}
\newcommand{\Bc}{\mathcal{B}}
\newcommand{\Cc}{\mathcal{C}}
\newcommand{\Ec}{\mathcal{E}}
\newcommand{\Pc}{\mathcal{P}}
\newcommand{\Nc}{\mathcal{N}}
\newcommand{\LN}{\mathcal{LN}}
\newcommand{\Expo}{\mathcal{E}}
\newcommand{\Poi}{\mathcal{P}}
\newcommand{\Ff}{\mathbb{F}}
\newcommand{\1}{\mathbf{1}}
\newcommand{\dd}{\,\mathrm{d}}
\DeclareMathOperator{\Var}{\mathbb{V}}
\DeclareMathOperator{\Cov}{Cov}
\DeclareMathOperator{\card}{card}
\DeclareMathOperator{\diag}{diag}
\newcommand{\indep}{\perp\!\!\!\perp}
\newcommand{\ps}{\ \PP\text{-p.s.}}
\newcommand{\tendloi}{\xrightarrow[n\to+\infty]{\text{loi}}}
\newcommand{\tendps}{\xrightarrow[n\to+\infty]{\text{p.s.}}}
\newcommand{\tendproba}{\xrightarrow[n\to+\infty]{\PP}}
\newcommand{\tendLp}{\xrightarrow[n\to+\infty]{L^p}}
\newcommand{\tend}{\xrightarrow[n\to+\infty]{}}

%     \fiche{C5}{Vecteurs gaussiens}{Gaussian vectors}
% ==============================================================
\usepackage[T1]{fontenc}
\usepackage[utf8]{inputenc}
\usepackage[french]{babel}
\usepackage{lmodern}
\usepackage[a4paper,margin=1.9cm,top=1.7cm,bottom=1.9cm]{geometry}
\usepackage{amsmath,amssymb,amsthm,mathtools}
\usepackage[inline]{enumitem}
\usepackage{booktabs,array,tabularx}
\usepackage{xcolor}
\usepackage[most]{tcolorbox}
\usepackage{fancyhdr}
\usepackage{hyperref}
\hypersetup{colorlinks=true,linkcolor=blue!50!black,urlcolor=blue!50!black}

\setlist{itemsep=1pt,topsep=3pt,parsep=0pt}
\setlength{\parindent}{0pt}
\setlength{\parskip}{3pt}
\emergencystretch=3em

% ---------- Identification de la fiche -------------------------
\newcommand{\ficheid}{}
\newcommand{\fichetitre}{}
\newcommand{\fiche}[3]{%
  \renewcommand{\ficheid}{#1}\renewcommand{\fichetitre}{#2}%
  \pagestyle{fancy}\fancyhf{}%
  \fancyhead[L]{\small\textsf{Advanced Probabilities -- Fiche #1}}%
  \fancyhead[R]{\small\textsf{#2}}%
  \fancyfoot[C]{\small\thepage}%
  \begin{center}
    {\LARGE\bfseries Fiche #1 -- #2}\par\smallskip
    {\large\itshape (#3)}\par\smallskip
    {\small Cours de V.\,M.~Lambert, M1 -- semestre 1}
  \end{center}\vspace{-2pt}\hrule\vspace{6pt}}

% ---------- Environnements numérotés (un seul compteur) --------
\theoremstyle{definition}
\newtheorem{definition}{Définition}
\renewcommand{\thedefinition}{\ficheid.\arabic{definition}}
\newtheorem{theoreme}[definition]{Théorème}
\newtheorem{proposition}[definition]{Proposition}
\newtheorem{lemme}[definition]{Lemme}
\newtheorem{corollaire}[definition]{Corollaire}
\newtheorem{exemple}[definition]{Exemple}
\theoremstyle{remark}
\newtheorem{remarque}[definition]{Remarque}

% Résultat hors cours (standard, signalé comme tel)
\theoremstyle{definition}
\newtheorem{horscours}[definition]{Hors cours}

% Idée de preuve
\newenvironment{idee}{\par\small\textit{Idée de preuve.}~}{\par}

% ---------- Encadrés -----------------------------------------
\tcbset{fichebox/.style={enhanced,breakable,boxrule=0.6pt,arc=2pt,
  left=5pt,right=5pt,top=3pt,bottom=3pt,fonttitle=\bfseries}}
\newtcolorbox{piege}[1][]{fichebox,colback=red!4,colframe=red!60!black,
  title={Pièges classiques},#1}
\newtcolorbox{liens}[1][]{fichebox,colback=blue!3,colframe=blue!50!black,
  title={Liens entre notions},#1}
\newtcolorbox{methode}[1][Méthode]{fichebox,colback=green!4,colframe=green!40!black,title={#1}}
\newtcolorbox{alerte}[1][]{fichebox,colback=orange!6,colframe=orange!70!black,
  title={Attention à l'énoncé},#1}

% Exercice de TD (fiches de rappels)
\newcounter{exo}
\newenvironment{exo}[1][]{\refstepcounter{exo}\par\medskip
  \noindent{\large\bfseries Exercice~\theexo}\ifx\relax#1\relax\else~{\large\bfseries -- #1}\fi\par\nopagebreak\smallskip}{\par}
\newcommand{\outils}{\textbf{Outils du cours.}~}
\newcommand{\meth}{\textbf{Méthode / astuce.}~}
\newcommand{\indic}{\textbf{Indication.}~}
\newcommand{\rep}{\textbf{Réponse.}~}

% ---------- Notations du cours ---------------------------------
\newcommand{\R}{\mathbb{R}}
\newcommand{\N}{\mathbb{N}}
\newcommand{\Q}{\mathbb{Q}}
\newcommand{\C}{\mathbb{C}}
\newcommand{\PP}{\mathbb{P}}
\newcommand{\QQ}{\mathbb{Q}}
\newcommand{\E}{\mathbb{E}}
\newcommand{\F}{\mathcal{F}}
\newcommand{\G}{\mathcal{G}}
\newcommand{\Hc}{\mathcal{H}}
\newcommand{\Tc}{\mathcal{T}}
\newcommand{\Bc}{\mathcal{B}}
\newcommand{\Cc}{\mathcal{C}}
\newcommand{\Ec}{\mathcal{E}}
\newcommand{\Pc}{\mathcal{P}}
\newcommand{\Nc}{\mathcal{N}}
\newcommand{\LN}{\mathcal{LN}}
\newcommand{\Expo}{\mathcal{E}}
\newcommand{\Poi}{\mathcal{P}}
\newcommand{\Ff}{\mathbb{F}}
\newcommand{\1}{\mathbf{1}}
\newcommand{\dd}{\,\mathrm{d}}
\DeclareMathOperator{\Var}{\mathbb{V}}
\DeclareMathOperator{\Cov}{Cov}
\DeclareMathOperator{\card}{card}
\DeclareMathOperator{\diag}{diag}
\newcommand{\indep}{\perp\!\!\!\perp}
\newcommand{\ps}{\ \PP\text{-p.s.}}
\newcommand{\tendloi}{\xrightarrow[n\to+\infty]{\text{loi}}}
\newcommand{\tendps}{\xrightarrow[n\to+\infty]{\text{p.s.}}}
\newcommand{\tendproba}{\xrightarrow[n\to+\infty]{\PP}}
\newcommand{\tendLp}{\xrightarrow[n\to+\infty]{L^p}}
\newcommand{\tend}{\xrightarrow[n\to+\infty]{}}

%     \fiche{C5}{Vecteurs gaussiens}{Gaussian vectors}
% ==============================================================
\usepackage[T1]{fontenc}
\usepackage[utf8]{inputenc}
\usepackage[french]{babel}
\usepackage{lmodern}
\usepackage[a4paper,margin=1.9cm,top=1.7cm,bottom=1.9cm]{geometry}
\usepackage{amsmath,amssymb,amsthm,mathtools}
\usepackage[inline]{enumitem}
\usepackage{booktabs,array,tabularx}
\usepackage{xcolor}
\usepackage[most]{tcolorbox}
\usepackage{fancyhdr}
\usepackage{hyperref}
\hypersetup{colorlinks=true,linkcolor=blue!50!black,urlcolor=blue!50!black}

\setlist{itemsep=1pt,topsep=3pt,parsep=0pt}
\setlength{\parindent}{0pt}
\setlength{\parskip}{3pt}
\emergencystretch=3em

% ---------- Identification de la fiche -------------------------
\newcommand{\ficheid}{}
\newcommand{\fichetitre}{}
\newcommand{\fiche}[3]{%
  \renewcommand{\ficheid}{#1}\renewcommand{\fichetitre}{#2}%
  \pagestyle{fancy}\fancyhf{}%
  \fancyhead[L]{\small\textsf{Advanced Probabilities -- Fiche #1}}%
  \fancyhead[R]{\small\textsf{#2}}%
  \fancyfoot[C]{\small\thepage}%
  \begin{center}
    {\LARGE\bfseries Fiche #1 -- #2}\par\smallskip
    {\large\itshape (#3)}\par\smallskip
    {\small Cours de V.\,M.~Lambert, M1 -- semestre 1}
  \end{center}\vspace{-2pt}\hrule\vspace{6pt}}

% ---------- Environnements numérotés (un seul compteur) --------
\theoremstyle{definition}
\newtheorem{definition}{Définition}
\renewcommand{\thedefinition}{\ficheid.\arabic{definition}}
\newtheorem{theoreme}[definition]{Théorème}
\newtheorem{proposition}[definition]{Proposition}
\newtheorem{lemme}[definition]{Lemme}
\newtheorem{corollaire}[definition]{Corollaire}
\newtheorem{exemple}[definition]{Exemple}
\theoremstyle{remark}
\newtheorem{remarque}[definition]{Remarque}

% Résultat hors cours (standard, signalé comme tel)
\theoremstyle{definition}
\newtheorem{horscours}[definition]{Hors cours}

% Idée de preuve
\newenvironment{idee}{\par\small\textit{Idée de preuve.}~}{\par}

% ---------- Encadrés -----------------------------------------
\tcbset{fichebox/.style={enhanced,breakable,boxrule=0.6pt,arc=2pt,
  left=5pt,right=5pt,top=3pt,bottom=3pt,fonttitle=\bfseries}}
\newtcolorbox{piege}[1][]{fichebox,colback=red!4,colframe=red!60!black,
  title={Pièges classiques},#1}
\newtcolorbox{liens}[1][]{fichebox,colback=blue!3,colframe=blue!50!black,
  title={Liens entre notions},#1}
\newtcolorbox{methode}[1][Méthode]{fichebox,colback=green!4,colframe=green!40!black,title={#1}}
\newtcolorbox{alerte}[1][]{fichebox,colback=orange!6,colframe=orange!70!black,
  title={Attention à l'énoncé},#1}

% Exercice de TD (fiches de rappels)
\newcounter{exo}
\newenvironment{exo}[1][]{\refstepcounter{exo}\par\medskip
  \noindent{\large\bfseries Exercice~\theexo}\ifx\relax#1\relax\else~{\large\bfseries -- #1}\fi\par\nopagebreak\smallskip}{\par}
\newcommand{\outils}{\textbf{Outils du cours.}~}
\newcommand{\meth}{\textbf{Méthode / astuce.}~}
\newcommand{\indic}{\textbf{Indication.}~}
\newcommand{\rep}{\textbf{Réponse.}~}

% ---------- Notations du cours ---------------------------------
\newcommand{\R}{\mathbb{R}}
\newcommand{\N}{\mathbb{N}}
\newcommand{\Q}{\mathbb{Q}}
\newcommand{\C}{\mathbb{C}}
\newcommand{\PP}{\mathbb{P}}
\newcommand{\QQ}{\mathbb{Q}}
\newcommand{\E}{\mathbb{E}}
\newcommand{\F}{\mathcal{F}}
\newcommand{\G}{\mathcal{G}}
\newcommand{\Hc}{\mathcal{H}}
\newcommand{\Tc}{\mathcal{T}}
\newcommand{\Bc}{\mathcal{B}}
\newcommand{\Cc}{\mathcal{C}}
\newcommand{\Ec}{\mathcal{E}}
\newcommand{\Pc}{\mathcal{P}}
\newcommand{\Nc}{\mathcal{N}}
\newcommand{\LN}{\mathcal{LN}}
\newcommand{\Expo}{\mathcal{E}}
\newcommand{\Poi}{\mathcal{P}}
\newcommand{\Ff}{\mathbb{F}}
\newcommand{\1}{\mathbf{1}}
\newcommand{\dd}{\,\mathrm{d}}
\DeclareMathOperator{\Var}{\mathbb{V}}
\DeclareMathOperator{\Cov}{Cov}
\DeclareMathOperator{\card}{card}
\DeclareMathOperator{\diag}{diag}
\newcommand{\indep}{\perp\!\!\!\perp}
\newcommand{\ps}{\ \PP\text{-p.s.}}
\newcommand{\tendloi}{\xrightarrow[n\to+\infty]{\text{loi}}}
\newcommand{\tendps}{\xrightarrow[n\to+\infty]{\text{p.s.}}}
\newcommand{\tendproba}{\xrightarrow[n\to+\infty]{\PP}}
\newcommand{\tendLp}{\xrightarrow[n\to+\infty]{L^p}}
\newcommand{\tend}{\xrightarrow[n\to+\infty]{}}


% ---------- Macros locales au corrigé --------------------------
\newtcolorbox{resultat}{fichebox,colback=blue!4,colframe=blue!50!black}
\newcommand{\question}[2]{\par\medskip\noindent\textbf{\ifx\relax#1\relax Question.\else Question~#1.\fi}~\textit{#2}\par\smallskip}
\providecommand{\Z}{\mathbb{Z}}
\newcommand{\Geo}{\mathcal{G}}
\newcommand{\Mult}{\mathcal{M}}
\newcommand{\Bin}{\mathcal{B}}

\begin{document}
\fiche{TD4}{Corrigé -- Espérance conditionnelle}{Conditional Expectation}

\textit{Corrigé détaillé du TD4. Renvois : « C6.5, point 7 » = fiche C6, énoncé n°5, item 7. Fiche de rappels associée :
\texttt{fiches/td/TD4\_esperance\_conditionnelle}. Sauf mention contraire, les paramètres de Poisson sont $>0$.}

% =================================================================
\begin{exo}[Somme de variables de Poisson indépendantes]
\question{}{$(X_n)$ indépendantes, $X_n\sim\Poi(\lambda_n)$. Montrer que $\sum_{n=1}^NX_n\sim\Poi\bigl(\sum_{n=1}^N\lambda_n\bigr)$ pour tout $N\ge1$.}

\textbf{Étape 1 : deux variables.} Soient $U\sim\Poi(a)$ et $V\sim\Poi(b)$ indépendantes. $U+V$ est à valeurs dans $\N$ et, pour $k\in\N$,
l'événement $\{U+V=k\}$ est la réunion \emph{disjointe} des $\{U=j,\,V=k-j\}$, $0\le j\le k$ (car $U,V\ge0$). Par additivité puis
indépendance (C3.8, avec $A_1=\{j\}$, $A_2=\{k-j\}$) :
\begin{align*}\PP(U+V=k)&=\sum_{j=0}^k\PP(U=j)\PP(V=k-j)=\sum_{j=0}^ke^{-a}\frac{a^j}{j!}\,e^{-b}\frac{b^{k-j}}{(k-j)!}\\
&=\frac{e^{-(a+b)}}{k!}\sum_{j=0}^k\binom kja^jb^{k-j}=e^{-(a+b)}\frac{(a+b)^k}{k!}\end{align*}
(binôme de Newton). Donc $U+V\sim\Poi(a+b)$ (C2.11).

\textbf{Étape 2 : récurrence sur $N$.} Notons $S_N=\sum_{n=1}^NX_n$ et $\Lambda_N=\sum_{n=1}^N\lambda_n$. Pour $N=1$ c'est l'hypothèse.
Si $S_N\sim\Poi(\Lambda_N)$, alors $S_N$ est une fonction mesurable du paquet $(X_1,\dots,X_N)$, disjoint de $\{X_{N+1}\}$ : par
regroupement par paquets (C3.11, hors cours), $S_N\indep X_{N+1}$. L'étape 1 avec $U=S_N$, $V=X_{N+1}$ donne
$S_{N+1}=S_N+X_{N+1}\sim\Poi(\Lambda_N+\lambda_{N+1})=\Poi(\Lambda_{N+1})$.

\begin{resultat}
Pour tout $N\ge1$ : $\displaystyle\sum_{n=1}^NX_n\sim\Poi\Bigl(\sum_{n=1}^N\lambda_n\Bigr)$.
\end{resultat}

\begin{remarque}[variantes]
\emph{(a) Fonctions caractéristiques.} Par C4.2, point 3 (et récurrence avec C3.11, hors cours), $\varphi_{S_N}=\prod_n\varphi_{X_n}$ ; avec
$\varphi_{X_n}(t)=e^{\lambda_n(e^{it}-1)}$ (C4.2, point 4, formule hors cours), $\varphi_{S_N}(t)=e^{\Lambda_N(e^{it}-1)}$, fonction caractéristique
de $\Poi(\Lambda_N)$, et C4.2, point 1 conclut.
\emph{(b) Fonction génératrice.} Pour $u\in[0,1]$, C3.14 (point 2, $h_n(x)=u^x\ge0$) donne $\E(u^{S_N})=\prod_n e^{\lambda_n(u-1)}=e^{\Lambda_N(u-1)}$
(C2.11). Hors cours : la fonction génératrice caractérise la loi d'une v.a.\ à valeurs dans $\N$ (unicité des coefficients d'une série entière).
\end{remarque}
\end{exo}

% =================================================================
\begin{exo}[Loi conditionnelle multinomiale]
Notations : $S=X_1+\dots+X_p$, $\Lambda=\lambda_1+\dots+\lambda_p>0$. Par l'exercice~1, $S\sim\Poi(\Lambda)$, donc
$\PP(S=n)=e^{-\Lambda}\Lambda^n/n!>0$ pour tout $n\in\N$ : le conditionnement par $\{S=n\}$ est élémentaire (C1.12).

\question{1}{Calculer $\PP(X_1=n_1,\dots,X_{p-1}=n_{p-1}\mid S=n)$ et reconnaître la loi de $(X_1,\dots,X_{p-1})$ sachant $S=n$.}

Posons $n_p:=n-\sum_{i=1}^{p-1}n_i\in\Z$.
\begin{itemize}
  \item Si $n_p<0$ : sur $\{X_1=n_1,\dots,X_{p-1}=n_{p-1},S=n\}$ on aurait $X_p=n_p<0$, impossible ; l'événement est vide et la probabilité
  conditionnelle vaut $0$.
  \item Si $n_p\ge0$ : $\{X_1=n_1,\dots,X_{p-1}=n_{p-1},S=n\}=\{X_1=n_1,\dots,X_{p-1}=n_{p-1},X_p=n_p\}$. Par indépendance (C3.8) :
  \[\PP(X_1=n_1,\dots,X_p=n_p)=\prod_{i=1}^pe^{-\lambda_i}\frac{\lambda_i^{n_i}}{n_i!}=e^{-\Lambda}\prod_{i=1}^p\frac{\lambda_i^{n_i}}{n_i!}.\]
  Par C1.12, et comme $\Lambda^n=\prod_i\Lambda^{n_i}$ (car $\sum_in_i=n$) :
  \[\PP(X_1=n_1,\dots,X_{p-1}=n_{p-1}\mid S=n)=\frac{e^{-\Lambda}\prod_i\lambda_i^{n_i}/n_i!}{e^{-\Lambda}\Lambda^n/n!}
  =\frac{n!}{n_1!\cdots n_p!}\prod_{i=1}^p\Bigl(\frac{\lambda_i}{\Lambda}\Bigr)^{n_i}.\]
\end{itemize}

\begin{resultat}
Pour $n_1,\dots,n_{p-1}\in\N$ et $n_p=n-\sum_{i<p}n_i$ :
\[\PP(X_1=n_1,\dots,X_{p-1}=n_{p-1}\mid S=n)=\begin{cases}\dfrac{n!}{n_1!\cdots n_p!}\displaystyle\prod_{i=1}^p\Bigl(\frac{\lambda_i}{\Lambda}\Bigr)^{n_i}&\text{si }n_p\ge0,\\[2mm]
0&\text{sinon.}\end{cases}\]
Sachant $\{S=n\}$, $(X_1,\dots,X_p)$ suit la loi \textbf{multinomiale} $\Mult\bigl(n;\frac{\lambda_1}\Lambda,\dots,\frac{\lambda_p}\Lambda\bigr)$ ;
$(X_1,\dots,X_{p-1})$ en est la marginale des $p-1$ premières coordonnées (la dernière vaut $X_p=n-\sum_{i<p}X_i$).
\end{resultat}

Interprétation : on répartit $n$ boules indépendamment dans $p$ urnes, l'urne $i$ étant choisie avec probabilité $\lambda_i/\Lambda$.
Cas $p=2$ : sachant $X_1+X_2=n$, $X_1\sim\Bin\bigl(n,\frac{\lambda_1}{\lambda_1+\lambda_2}\bigr)$ (C2.11), puisque
$\frac{n!}{n_1!(n-n_1)!}q^{n_1}(1-q)^{n-n_1}$ avec $q=\frac{\lambda_1}{\lambda_1+\lambda_2}$.

\question{2}{Calculer $\E(X_1\mid X_1+X_2)$.}

Posons $T=X_1+X_2$ et $q=\frac{\lambda_1}{\lambda_1+\lambda_2}$. On a $X_1\in L^1$ ($\E X_1=\lambda_1$) et $T$ est discrète, à valeurs dans $\N$,
avec $\PP(T=m)>0$ pour tout $m$ (exercice~1, $T\sim\Poi(\lambda_1+\lambda_2)$). On applique C6.2 :
$\E(X_1\mid T)=\phi(T)$ avec, pour $m\in\N$,
\[\phi(m)=\frac{\E(X_1\1_{\{T=m\}})}{\PP(T=m)}=\sum_{j=0}^m j\,\frac{\PP(X_1=j,T=m)}{\PP(T=m)}=\sum_{j=0}^mj\,\PP(X_1=j\mid T=m)=mq,\]
car, d'après la question~1 (cas $p=2$), $\PP(X_1=\cdot\mid T=m)$ est la loi $\Bin(m,q)$, d'espérance $mq$ (C2.11).
(On a écrit $X_1\1_{\{T=m\}}=\sum_{j=0}^mj\1_{\{X_1=j,T=m\}}$, somme finie.)

\begin{resultat}
$\displaystyle\E(X_1\mid X_1+X_2)=\frac{\lambda_1}{\lambda_1+\lambda_2}\,(X_1+X_2)$ \quad p.s.\ (ici même partout, par la définition C6.2).
\end{resultat}

\textbf{Vérifications.} Espérance totale (C6.5, point 1) : $\E\bigl(q(X_1+X_2)\bigr)=q(\lambda_1+\lambda_2)=\lambda_1=\E X_1$.
De même $\E(X_2\mid T)=(1-q)T$, et la somme vaut $T=\E(T\mid T)$, conformément à C6.5, points 2 et 4.
\end{exo}

% =================================================================
\begin{exo}[Mélange Poisson--exponentielle]
\question{}{$X$ à valeurs dans $\N$, $Y\sim\Expo(1)$, et la loi de $X$ sachant $\{Y=y\}$ est $\Poi(y)$. Lois de $(X,Y)$, de $X$, et
de $Y$ sachant $\{X=k\}$.}

\textbf{Sens de l'hypothèse.} Comme $\PP(Y=y)=0$, « sachant $\{Y=y\}$ » ne relève pas de C1.12 ; par analogie avec C6.9 (loi jointe
$=$ loi conditionnelle $\times$ loi marginale, cas mixte non traité dans le cours), l'hypothèse signifie :
\begin{equation}\label{eq:td4-melange}
\forall k\in\N,\ \forall B\in\Bc(\R),\qquad \PP(X=k,\,Y\in B)=\int_B e^{-y}\frac{y^k}{k!}\;\underbrace{e^{-y}\1_{y>0}}_{f_Y(y)}\dd y .
\end{equation}
(Hors cours : notion de loi conditionnelle / noyau de transition.)

\textbf{Loi du couple.} Posons $g(k,y)=e^{-2y}\dfrac{y^k}{k!}\1_{y>0}$ sur $\N\times\R$. Par \eqref{eq:td4-melange},
$\PP\bigl((X,Y)\in\{k\}\times B\bigr)=\sum_{j\in\N}\int\1_{\{k\}\times B}(j,y)\,g(j,y)\dd y$ : la loi de $(X,Y)$ et la mesure de densité $g$
par rapport à (comptage $\otimes$ Lebesgue) (C2.6, C3.16 : mesure de référence produit) coïncident sur les pavés $\{k\}\times B$, qui forment un $\pi$-système
engendrant $\mathcal P(\N)\otimes\Bc(\R)$ ; elles sont donc égales (unicité, argument de classe monotone, hors cours).

\begin{resultat}
$(X,Y)$ a pour densité $g(k,y)=e^{-2y}\dfrac{y^k}{k!}\1_{y>0}$ par rapport à (comptage sur $\N$) $\otimes$ (Lebesgue sur $\R$) :
$\PP(X=k,\,Y\in B)=\displaystyle\int_B e^{-2y}\frac{y^k}{k!}\1_{y>0}\dd y$.
\end{resultat}

Cohérence : $\sum_{k\ge0}g(k,y)=e^{-2y}e^{y}\1_{y>0}=e^{-y}\1_{y>0}$, on retrouve bien $Y\sim\Expo(1)$ (sommation terme à terme d'une série positive).

\textbf{Loi de $X$.} Avec $B=\R$ dans \eqref{eq:td4-melange}, et en reconnaissant la densité d'Erlang $f_{k+1}$ de paramètre $\lambda=2$,
$f_{k+1}(y)=2^{k+1}e^{-2y}\frac{y^k}{k!}\1_{y>0}$, d'intégrale $1$ :
\[\PP(X=k)=\int_0^{\infty}e^{-2y}\frac{y^k}{k!}\dd y=\frac1{2^{k+1}}\int_\R f_{k+1}(y)\dd y=\frac1{2^{k+1}},\qquad k\in\N.\]
Ainsi $\PP(X+1=j)=\bigl(1-\tfrac12\bigr)^{j-1}\tfrac12$ pour $j\ge1$ : $X+1$ suit la loi géométrique sur $\N^*$ de paramètre $\frac12$ (C2.11).
(Contrôle numérique par intégration au point milieu sous \texttt{awk}, $k=0,\dots,4$ : $0{,}500$ ; $0{,}250$ ; $0{,}125$ ; $0{,}0625$ ; $0{,}0313$.)

\begin{resultat}
$\PP(X=k)=\dfrac1{2^{k+1}}$, $k\in\N$ : $X+1\sim\Geo(\tfrac12)$ (géométrique sur $\N^*$), $\E X=1$.
\end{resultat}

\textbf{Loi de $Y$ sachant $\{X=k\}$.} Ici $\PP(X=k)=2^{-(k+1)}>0$ : conditionnement élémentaire (C1.12). Pour $B\in\Bc(\R)$,
\[\PP(Y\in B\mid X=k)=\frac{\PP(X=k,Y\in B)}{\PP(X=k)}=\int_B2^{k+1}e^{-2y}\frac{y^k}{k!}\1_{y>0}\dd y=\int_Bf_{k+1}(y)\dd y .\]

\begin{resultat}
Sachant $\{X=k\}$, $Y$ a pour densité $2^{k+1}\dfrac{y^k}{k!}e^{-2y}\1_{y>0}$ : loi d'\textbf{Erlang d'ordre $k+1$ et de paramètre $2$}.
\end{resultat}

\begin{remarque}
En conséquence (C6.2, Erlang$(k+1,2)$ d'espérance $\frac{k+1}2$, C2.11) : $\E(Y\mid X)=\frac{X+1}2$ p.s. ; et
$\E\bigl(\frac{X+1}2\bigr)=1=\E Y$ (C6.5, point 1). Lecture bayésienne : loi a priori $\Expo(1)=$ Erlang$(1,1)$, observation $\Poi(y)$,
loi a posteriori Erlang$(k+1,2)$.
\end{remarque}
\end{exo}

% =================================================================
\begin{exo}[Maximum et minimum]
Notations : $X=\max_iX_i$, $Y=\min_iX_i$, $F(t)=\PP(X_1\le t)=\int_{-\infty}^tf(s)\dd s$. On prend $n\ge2$ (pour $n=1$, $X=Y=X_1$).

\textbf{Préliminaire : pas d'ex \ae quo.} Pour $i\ne j$, $(X_i,X_j)$ a pour densité $f(u)f(v)$ (C3.16, point 1), donc
$\PP(X_i=X_j)=\iint\1_{\{u=v\}}f(u)f(v)\dd u\dd v=0$ (la diagonale est Lebesgue-négligeable ; Fubini--Tonelli, hors cours).
Ainsi, p.s., les $X_i$ sont deux à deux distincts.

\question{1}{Calculer $F_X$, $F_Y$ et $F_{(X,Y)}(x,y)$ pour $x\le y$ et $x>y$.}

Par indépendance (C3.8, ou C3.9) :
\[F_X(x)=\PP\Bigl(\bigcap_{i=1}^n\{X_i\le x\}\Bigr)=\prod_{i=1}^n\PP(X_i\le x)=F(x)^n,\qquad
\PP(Y>y)=\PP\Bigl(\bigcap_{i=1}^n\{X_i>y\}\Bigr)=(1-F(y))^n .\]
Pour le couple, $F_{(X,Y)}(x,y)=\PP(X\le x,Y\le y)$.
\begin{itemize}
  \item Si $x\le y$ : $Y\le X$, donc $\{X\le x\}\subset\{Y\le y\}$ et $F_{(X,Y)}(x,y)=\PP(X\le x)=F(x)^n$.
  \item Si $x>y$ : $\PP(X\le x,Y\le y)=\PP(X\le x)-\PP(X\le x,Y>y)$ et
  $\{X\le x,Y>y\}=\bigcap_i\{y<X_i\le x\}$, de probabilité $\prod_i\PP(y<X_i\le x)=(F(x)-F(y))^n$ (C3.8).
\end{itemize}

\begin{resultat}
$F_X=F^n$, \quad $F_Y=1-(1-F)^n$, \quad
$F_{(X,Y)}(x,y)=\begin{cases}F(x)^n&\text{si }x\le y,\\F(x)^n-\bigl(F(x)-F(y)\bigr)^n&\text{si }x>y.\end{cases}$
\end{resultat}

\question{2}{Montrer que $f_X=nf\,F^{n-1}$ ; densité de $Y$.}

Méthode de la fonction muette (C2.13, C2.7). Soit $h\ge0$ mesurable. Notons $A_i=\{X_k<X_i\ \forall k\ne i\}$ (« le maximum est atteint
en $i$ ») : les $A_i$ sont disjoints et, d'après le préliminaire, $\PP\bigl(\bigcup_iA_i\bigr)=1$. Sur $A_i$, $X=X_i$, donc
$h(X)=\sum_ih(X_i)\1_{A_i}$ p.s. Le vecteur $(X_1,\dots,X_n)$ a pour densité $\prod_kf(x_k)$ (C3.16, point 1) ; par transfert (C2.8,
C2.9) et Fubini--Tonelli (hors cours) en intégrant d'abord en $x_k$, $k\ne i$ :
\[\E\bigl(h(X_i)\1_{A_i}\bigr)=\int_{\R^n}h(x_i)\prod_{k\ne i}\1_{\{x_k<x_i\}}\prod_{k=1}^nf(x_k)\dd x
=\int_\R h(t)f(t)F(t)^{n-1}\dd t .\]
En sommant sur $i$ : $\E(h(X))=\int_\R h(t)\,nf(t)F(t)^{n-1}\dd t$, ce qui identifie la densité de $X$. De même, avec
$A'_i=\{X_k>X_i\ \forall k\ne i\}$ et $\PP(X_k>t)=1-F(t)$ : $\E(h(Y))=\int_\R h(t)\,nf(t)(1-F(t))^{n-1}\dd t$.

\begin{resultat}
$f_X(x)=nf(x)F(x)^{n-1}$, \qquad $f_Y(y)=nf(y)\bigl(1-F(y)\bigr)^{n-1}$.
\end{resultat}
Cohérence avec la question 1 : aux points où $f$ est continue, $F'=f$ et $(F^n)'=nfF^{n-1}$, $(1-(1-F)^n)'=nf(1-F)^{n-1}$
(c'est l'alternative « dériver la fonction de répartition » de C2.13).

\question{3}{Montrer que $f_{(Y,X)}(y,x)=n(n-1)f(y)f(x)(F(x)-F(y))^{n-2}\1_{y<x}$.}

Même méthode. Pour $i\ne j$, soit $B_{ij}=\{X_i<X_k<X_j\ \forall k\notin\{i,j\}\}$ (« le min est en $i$, le max en $j$ »). Les $n(n-1)$
événements $B_{ij}$ sont disjoints et, p.s.\ (préliminaire), l'un d'eux est réalisé ; sur $B_{ij}$, $(Y,X)=(X_i,X_j)$ et $X_i<X_j$. Pour
$h:\R^2\to\R_+$ mesurable, avec C3.16, C2.8 et Fubini--Tonelli (hors cours), en intégrant d'abord les $n-2$ autres coordonnées :
\[\E\bigl(h(Y,X)\bigr)=\sum_{i\ne j}\E\bigl(h(X_i,X_j)\1_{B_{ij}}\bigr)
=n(n-1)\iint_{\R^2}h(y,x)\1_{\{y<x\}}f(y)f(x)\bigl(F(x)-F(y)\bigr)^{n-2}\dd y\dd x,\]
car $\int_\R\1_{\{y<t<x\}}f(t)\dd t=F(x)-F(y)$ pour chacune des $n-2$ coordonnées restantes (convention $0^0=1$ si $n=2$). D'où la densité
annoncée (C2.7).

\begin{resultat}
$f_{(Y,X)}(y,x)=n(n-1)f(y)f(x)\bigl(F(x)-F(y)\bigr)^{n-2}\1_{y<x}$.
\end{resultat}
Cohérence avec la question~1 : pour $x>y$, $\frac{\partial^2}{\partial x\,\partial y}\bigl[F(x)^n-(F(x)-F(y))^n\bigr]=n(n-1)f(x)f(y)(F(x)-F(y))^{n-2}$
(points de continuité de $f$), et la dérivée croisée est nulle pour $x\le y$.

\question{4}{Loi conditionnelle de $Y$ sachant $\{X=x\}$ et $\E(Y\mid X)$.}

On applique C6.9 au couple $(X,Y)$ (de densité $(x,y)\mapsto f_{(Y,X)}(y,x)$). La marginale $\int_\R f_{(Y,X)}(y,x)\dd y$ est une densité de $X$ ;
elle coïncide avec celle de la question 2 (Lebesgue-p.p.\ par unicité de la densité, et même partout puisque $y\mapsto-(F(x)-F(y))^{n-1}/(n-1)$ est une primitive de l'intégrande sur $]-\infty,x[$), et l'on peut prendre la version $f_X(x)=nf(x)F(x)^{n-1}$ (la loi
conditionnelle n'est définie que pour $\PP_X$-presque tout $x$). Pour $x$ tel que $f_X(x)>0$ (i.e.\ $f(x)>0$ et $F(x)>0$) :
\begin{align*}f_Y^{X=x}(y)&=\frac{f_{(Y,X)}(y,x)}{f_X(x)}=\frac{n(n-1)f(y)f(x)(F(x)-F(y))^{n-2}}{nf(x)F(x)^{n-1}}\1_{y<x}\\
&=(n-1)\frac{f(y)\bigl(F(x)-F(y)\bigr)^{n-2}}{F(x)^{n-1}}\1_{y<x}.\end{align*}
\textbf{Espérance conditionnelle.} Hypothèse : $X_1\in L^1$ ; alors $Y\in L^1$ car $|Y|\le\sum_i|X_i|$. C6.9 s'applique à $g(x,y)=y^+$
et $g(x,y)=y^-$ (positives) ; par linéarité (C6.5, point 2), $\E(Y\mid X)=\psi(X)$ p.s.\ avec
$\psi(x)=\int_\R y\,f_Y^{X=x}(y)\dd y$, fini pour $\PP_X$-presque tout $x$ (car $\int_\R\bigl(\int_\R|y|\,f_Y^{X=x}(y)\dd y\bigr)f_X(x)\dd x=\E|Y|<\infty$).

\begin{resultat}
Sachant $\{X=x\}$, $Y$ a pour densité $\displaystyle f_Y^{X=x}(y)=(n-1)\frac{f(y)\bigl(F(x)-F(y)\bigr)^{n-2}}{F(x)^{n-1}}\1_{y<x}$, et
\[\E(Y\mid X)=\psi(X)\ \text{p.s.},\qquad \psi(x)=\frac{n-1}{F(x)^{n-1}}\int_{-\infty}^{x}y\,f(y)\bigl(F(x)-F(y)\bigr)^{n-2}\dd y .\]
\end{resultat}
Cas $n=2$ : $f_Y^{X=x}(y)=\frac{f(y)}{F(x)}\1_{y<x}$, i.e.\ sachant que le max vaut $x$, l'autre variable suit la loi de $X_1$ conditionnée à
$\{X_1<x\}$ ; $\psi(x)=\E(X_1\1_{\{X_1<x\}})/F(x)$.

\question{5}{Cas uniforme sur $[0,1]$.}

Ici $f=\1_{[0,1]}$ et $F(t)=t$ sur $[0,1]$ ($0$ avant, $1$ après). Pour $x,y\in[0,1]$, les questions 1 à 3 donnent :
\[F_X(x)=x^n,\quad F_Y(y)=1-(1-y)^n,\quad F_{(X,Y)}(x,y)=\begin{cases}x^n&x\le y\\x^n-(x-y)^n&x>y\end{cases},\]
\[f_X(x)=nx^{n-1}\1_{[0,1]}(x),\quad f_Y(y)=n(1-y)^{n-1}\1_{[0,1]}(y),\quad f_{(Y,X)}(y,x)=n(n-1)(x-y)^{n-2}\1_{0\le y<x\le1}.\]
Pour $x\in\,]0,1]$ (on a bien $f_X(x)>0$) : $f_Y^{X=x}(y)=(n-1)\dfrac{(x-y)^{n-2}}{x^{n-1}}\1_{0\le y<x}$. Puis, par intégration par parties,
\[\int_0^xy\,(n-1)(x-y)^{n-2}\dd y=\Bigl[-y(x-y)^{n-1}\Bigr]_0^x+\int_0^x(x-y)^{n-1}\dd y=\frac{x^n}n,
\qquad\text{donc}\qquad\psi(x)=\frac{x^n/n}{x^{n-1}}=\frac xn .\]
(Contrôle \texttt{awk} par point milieu : $n=3$, $x=0{,}6$ : $\psi=0{,}200$ ; $n=5$, $x=0{,}9$ : $\psi=0{,}180$.)

\begin{resultat}
Uniforme : sachant $\{X=x\}$ ($0<x\le1$), $Y$ a pour densité $(n-1)\dfrac{(x-y)^{n-2}}{x^{n-1}}\1_{0\le y<x}$ (uniforme sur $[0,x]$ si $n=2$), et
\[\E(Y\mid X)=\frac Xn\quad\text{p.s.}\]
\end{resultat}
\textbf{Vérification} (C6.5, point 1) : $\E\bigl(\frac Xn\bigr)=\frac1n\int_0^1x\,nx^{n-1}\dd x=\frac1{n+1}$ et
$\E Y=\int_0^1y\,n(1-y)^{n-1}\dd y=\int_0^1(1-y)^n\dd y=\frac1{n+1}$ : cohérent.
\end{exo}

% =================================================================
\begin{exo}[Régression gaussienne]
Le vecteur $(X,Y_1,Y_2)$ est gaussien centré, de covariance $K=\begin{pmatrix}4&1&2\\1&9&-3\\2&-3&4\end{pmatrix}$. On lit
$\Var X=4$, $C=\begin{pmatrix}\Cov(X,Y_1)\\\Cov(X,Y_2)\end{pmatrix}=\begin{pmatrix}1\\2\end{pmatrix}$,
$\Sigma_Y=\begin{pmatrix}9&-3\\-3&4\end{pmatrix}$.
$K$ est bien une matrice de covariance : symétrique, de mineurs principaux dominants $4$, $35$, $\det K=56$, tous $>0$ (définie positive par le critère de Sylvester, hors cours).

\question{1}{Calculer $\E(X\mid Y_1,Y_2)$.}

Hypothèses de C6.10 : $(Y_1,Y_2,X)$ est un vecteur gaussien (image de $(X,Y_1,Y_2)$ par une permutation des coordonnées, C5.5) et
$\det\Sigma_Y=36-9=27\ne0$. Donc, avec $\E X=\E Y_j=0$,
\[\Sigma_Y^{-1}=\frac1{27}\begin{pmatrix}4&3\\3&9\end{pmatrix},\qquad
\Lambda=\Sigma_Y^{-1}C=\frac1{27}\begin{pmatrix}4+6\\3+18\end{pmatrix}=\begin{pmatrix}10/27\\21/27\end{pmatrix},\qquad
\E(X\mid Y_1,Y_2)=\Lambda^tY\ \text{p.s.}\]

\begin{resultat}
$\displaystyle\E(X\mid Y_1,Y_2)=\frac{10Y_1+21Y_2}{27}=\frac{10}{27}Y_1+\frac79Y_2\approx0{,}370\,Y_1+0{,}778\,Y_2$ \quad p.s.
\end{resultat}

\question{2}{Vérifier que $X-\E(X\mid Y_1,Y_2)\perp(Y_1,Y_2)$.}

Soit $R=X-\frac{10Y_1+21Y_2}{27}$. Par bilinéarité de la covariance :
\[\Cov(R,Y_1)=1-\frac{10\cdot9+21\cdot(-3)}{27}=1-\frac{27}{27}=0,\qquad
\Cov(R,Y_2)=2-\frac{10\cdot(-3)+21\cdot4}{27}=2-\frac{54}{27}=0 .\]
Tout est centré, donc $\E(RY_1)=\E(RY_2)=0$ : $R$ est \textbf{orthogonal} dans $L^2$ à $Y_1$, $Y_2$ (et à toute combinaison linéaire). Mieux :
$(R,Y_1,Y_2)$ est l'image de $(X,Y_1,Y_2)$ par une application linéaire, donc un vecteur gaussien (C5.5) ; ses blocs $R$ et $(Y_1,Y_2)$
étant non corrélés, ils sont \textbf{indépendants} (C5.12, $2\Rightarrow1$).

\begin{resultat}
$\Cov(R,Y_1)=\Cov(R,Y_2)=0$ : $R=X-\E(X\mid Y_1,Y_2)$ est orthogonal à $(Y_1,Y_2)$ dans $L^2$, et même indépendant de $(Y_1,Y_2)$.
\end{resultat}

\question{3}{Calculer $\E\bigl((X-\E(X\mid Y_1,Y_2))^2\mid Y_1,Y_2\bigr)$.}

$R^2$ est intégrable ($R$ gaussienne, C5.3) et $\sigma(R^2)\subset\sigma(R)$ est indépendante de $\sigma(Y_1,Y_2)$ (question 2). Par C6.5,
point 7 (indépendance) : $\E(R^2\mid Y_1,Y_2)=\E(R^2)=\Var R$ (car $\E R=0$). Or, avec $\Sigma_Y\Lambda=C$,
\[\Var R=\Var X-2\Lambda^tC+\Lambda^t\Sigma_Y\Lambda=\Var X-\Lambda^tC=4-\frac{10\cdot1+21\cdot2}{27}=4-\frac{52}{27}=\frac{56}{27}.\]
(Contrôle : $\det K/\det\Sigma_Y=56/27$, formule du complément de Schur, hors cours.)

\begin{resultat}
$\E\bigl((X-\E(X\mid Y_1,Y_2))^2\mid Y_1,Y_2\bigr)=\dfrac{56}{27}\approx2{,}07$ \quad p.s.\ (une constante).
\end{resultat}

\question{4}{Loi de $X$ sachant $(Y_1,Y_2)$.}

Écrivons $X=M+R$ avec $M=\frac{10Y_1+21Y_2}{27}$, réelle et $\sigma(Y_1,Y_2)$-mesurable, et $R\sim\Nc\bigl(0,\frac{56}{27}\bigr)$ (combinaison
linéaire d'un vecteur gaussien centré, C5.3, variance de la question~3), indépendante de $\G=\sigma(Y_1,Y_2)$ (question 2). Pour
$h:\R\to\R$ mesurable bornée, le lemme de gel (C6.6, avec $\Psi(m,r)=h(m+r)$ bornée) donne
\[\E\bigl(h(X)\mid Y_1,Y_2\bigr)=\varphi(M)\ \text{p.s.},\qquad \varphi(m)=\E\,h(m+R)=\int_\R h(t)\,\frac{1}{\sqrt{2\pi s^2}}e^{-\frac{(t-m)^2}{2s^2}}\dd t,\quad s^2=\frac{56}{27},\]
i.e.\ $\varphi(m)$ est l'intégrale de $h$ contre la loi $\Nc(m,s^2)$. C'est exactement C6.11 (hors cours).

\begin{resultat}
Sachant $(Y_1,Y_2)$, $X\sim\Nc\Bigl(\dfrac{10Y_1+21Y_2}{27},\ \dfrac{56}{27}\Bigr)$ ; écart-type conditionnel $\sqrt{56/27}\approx1{,}44$.
En particulier $\PP(X\le t\mid Y_1,Y_2)=\Phi\Bigl(\dfrac{t-(10Y_1+21Y_2)/27}{\sqrt{56/27}}\Bigr)$ p.s.
\end{resultat}
\textbf{Remarque.} La variance conditionnelle ne dépend pas de $(Y_1,Y_2)$ et $\frac{56}{27}<4=\Var X$ : l'observation de $(Y_1,Y_2)$ réduit
l'incertitude sur $X$ de $\frac{52}{27}$ (part de variance « expliquée » $\frac{52}{108}\approx48{,}1\,\%$).
\end{exo}

\end{document}
